% 六模式族 V 的逐坐标邻居（伪立方体）。两种视图描述同一组模式。
% 左：标签网格。同列两点只差 x1（蓝实线），同行两点只差 x2（绿虚线）。
%     每个实心点恰有一条蓝线、一条绿线，即在两个坐标上都有邻居。
% 右：同一关系画成六边形环：
%   h1(正常,正常) -x2- h2(正常,广告) -x1- h3(广告,广告) -x2- h4(广告,诈骗)
%   -x1- h5(诈骗,诈骗) -x2- h6(诈骗,正常) -x1- h1。
\begingroup
\definecolor{LtSixGreen}{HTML}{829C8B}
\begin{tikzpicture}[figure picture,x=1mm,y=1mm]
 \def\cs{11}
 \foreach \r in {1,2,3}{\foreach \c in {1,2,3}{
   \draw[FigureGrid,line width=.5pt] ({8+\cs*(\c-1)},{22-\cs*(\r-1)}) rectangle ++(\cs,\cs);}}
 \foreach \lab [count=\c] in {正常,广告,诈骗}{
   \node[anchor=south,font=\figurefont\scriptsize] at ({8+\cs*(\c-.5)},33.5) {\lab};}
 \foreach \lab [count=\r] in {正常,广告,诈骗}{
   \node[anchor=east,font=\figurefont\scriptsize] at (7,{27.5-\cs*(\r-1)}) {\lab};}
 \node[figure note,anchor=south] at ({8+1.5*\cs},38) {列：$x_2$；行：$x_1$};
 \foreach \r/\c/\n in {1/1/a,1/2/b,2/2/c,2/3/d,3/3/e,3/1/f}{
   \coordinate (g\n) at ({8+\cs*(\c-.5)},{27.5-\cs*(\r-1)});}
 % x1 changes (same column): blue solid
 \draw[FigureBlue,line width=1pt] (gb)--(gc);
 \draw[FigureBlue,line width=1pt] (gd)--(ge);
 \draw[FigureBlue,line width=1pt] (ga) to[bend right=35] (gf);
 % x2 changes (same row): green dashed
 \draw[LtSixGreen,line width=1pt,dash pattern=on 2.2pt off 1.4pt] (ga)--(gb);
 \draw[LtSixGreen,line width=1pt,dash pattern=on 2.2pt off 1.4pt] (gc)--(gd);
 \draw[LtSixGreen,line width=1pt,dash pattern=on 2.2pt off 1.4pt] (gf) to[bend right=30] (ge);
 \foreach \r/\c in {1/3,2/1,3/2}{
   \draw[FigureMuted,line width=.6pt,fill=white] ({8+\cs*(\c-.5)},{27.5-\cs*(\r-1)}) circle (1.2);}
 \foreach \n in {a,b,c,d,e,f}{\fill[FigureBlue] (g\n) circle (1.7);}
 % ---- right: hexagon
 \tikzset{ltsixv/.style={draw=FigureBlue,fill=FigureBlueFill,rounded corners=1.8pt,
   line width=.7pt,inner sep=2pt,minimum height=6mm,font=\figurefont\scriptsize}}
 \node[ltsixv] (h1) at (66,32) {$h_1$ 正常，正常};
 \node[ltsixv] (h2) at (92,32) {$h_2$ 正常，广告};
 \node[ltsixv] (h3) at (101,16.5) {$h_3$ 广告，广告};
 \node[ltsixv] (h4) at (92,1) {$h_4$ 广告，诈骗};
 \node[ltsixv] (h5) at (66,1) {$h_5$ 诈骗，诈骗};
 \node[ltsixv] (h6) at (57,16.5) {$h_6$ 诈骗，正常};
 \foreach \u/\v in {h1/h2,h3/h4,h5/h6}{
   \draw[LtSixGreen,line width=1pt,dash pattern=on 2.2pt off 1.4pt] (\u)--(\v);}
 \foreach \u/\v in {h2/h3,h4/h5,h6/h1}{
   \draw[FigureBlue,line width=1pt] (\u)--(\v);}
 % legend
 \draw[FigureBlue,line width=1pt] (8,-8)--(15,-8);
 \node[figure note,anchor=west] at (16,-8) {只改 $x_1$ 的邻居};
 \draw[LtSixGreen,line width=1pt,dash pattern=on 2.2pt off 1.4pt] (45,-8)--(52,-8);
 \node[figure note,anchor=west] at (53,-8) {只改 $x_2$ 的邻居};
\end{tikzpicture}
\endgroup
